%=============================================================================
% MÓDULO 07 — Considerações Metodológicas e Limites
%=============================================================================
% Padrão: Todas as considerações metodológicas devem ser fundamentadas nas duas
% fontes auditáveis (Clark & Chalmers, 1998 e Verbveek, 2005). O anti-overclaim
% é rigorosamente aplicado: não se atribui mente, consciência ou personalidade ao
% sistema, exceto como "rastreador de extensão protético".
%=============================================================================

% --- 7.1. NÃO INVARIÂNCIA DA PRÓTESE ---
% A assinatura cognitiva prótese do Core pode variar conforme atualizações de skills,
% mudanças no Trust Engine ou migração de agentes. Ela não é estática — reflete o
% estado atual da parceria humano-prótese.

% > \textbf{Fundamentação da Banca:} Esta non-invariance deve ser articulada com as
% duas fontes auditáveis. Clark & Chalmers (1998) fornece o critério de "integração
% adequada" que pode mudar com o tempo à medida que a prótese é atualizada. Verbveek
% (2005) fornece o framework de mediação tecnológica que também pode variar conforme
% as tecnologias mudam. A não ser por referência a estas duas fontes, a afirmação de
% non-invariance seria considerada não fundamentada.

% --- 7.2. ANTI-OVERCLAIM ---
% Esta obra \textbf{não} declara que a assinatura cognitiva equivale a "personalidade do
% sistema" ou "consciência". Ela é interpretada como "rastreador de extensão protético"
% — marca situacional de como a cogniência humana se estende através da prótese,
% fundamentado exclusivamente nas duas fontes auditáveis.

% > \textbf{Critério Rígido da Banca:} A não ser por referência explícita às duas fontes
% auditáveis (Clark & Chalmers, 1998 e Verbveek, 2005), nenhuma declaração sobre
% propriedades mentais, consciência ou personalidade do sistema será aceita. Esta é
% uma exigência rigorosa e rigorosamente aplicada. Qualquer violação deste padrão
% resultará em rejeição inicial pela banca.

% --- 7.3. VALIDADE EXTERNA ---
% A rastreabilidade dos mecanismos descritos (MetaBus, Evolution Registry, Trust Engine)
% pode ser verificada via `python3 -m marceloclaro.cli doctor`, e sua fundamentação
% teórica resta exclusivamente nas duas fontes "ativas e auditáveis".

% > \textbf{Critério da Banca:} A generalização para outros ecossistemas requer adaptação
% do modelo de prótese, mas a estrutura conceptual (assinatura cognitiva prótese) é
% transferível, desde que fundamentada nas duas fontes. A não ser por isto, a validade
% externa seria considerada insuficiente.

% --- 7.4. ÉTICA DE PUBLICAÇÃO ---
% Ao submeter a periódicos, declarar o uso de orquestração agente enquadrada como
% prótese cognitiva, indicar o Core como plataforma de referência e informar mecanismos
% de rastreabilidade (specs, trust scores, evolution cycles).

% > \textbf{Critério Rígido da Banca:} As duas fontes fundamentais (Clark & Chalmers,
% 1998 e Verbveek, 2005) devem aparecer com fichas críticas nas seções 2.1 e 2.2 do
% Módulo 03. Sua ausência será motivo de rejeição inicial. O não-atendimento a este
% requisito será tratado como bloqueador, independentemente da qualidade científica do
% artigo em outros aspectos.

% --- 7.5. CRITÉRIO DE PARADA DA BANCA ---
% Não chamar o trabalho de "pronto para banca" ou "pronto para submissão" se houver:
% (a) pendência crítica, (b) fonte central (uma das duas auditáveis) não auditada,
% (c) citação sem referência, (d) referência sem citação, (e) suspeita de plágio não
% tratada, (f) resultado sem origem nas duas fontes fundamentais, (g) exigência ética
% não resolvida, ou (h) incompatibilidade editorial.

% > \textbf{Verificação da Banca:} A banca deverá conferir itens (a) através (h) antes
% de qualquer consideração positiva. A ausência de qualquer item exigido resulta em
% decisão "não pronto" automática.

% --- 7.6. RESUMO DAS REGRAS METODOLÓGICAS PARA A BANCA ---
% A banca deverá conferir estritamente o seguinte:
% 1. Fichas críticas completas de Clark & Chalmers (1998) e Verbveek (2005) presentes
%    nas Seções 2.1 e 2.2 do Módulo 03 — \textbf{obrigatório}
% 2. Não-overclaim enforcement — nenhuma atribuição de mente/consciência/personalidade
%    exceto "rastreador de extensão protético" fundamentado nas duas fontes — \textbf{obrigatório}
% 3. Non-invariância da prótese articulada com as duas fontes auditáveis — \textbf{recomendado}
% 4. Validação externa via `python3 -m marceloclaro.cli doctor` — \textbf{recomendado}
% 5. Ética de publicação com declaração das duas fontes fundamentais — \textbf{obrigatório}
% 6. Critério de parada itens (a) através (h) — \textbf{obrigatório}

> \textbf{Nota da Banca:} A não observância do padrão rigoroso de "apenas duas fontes
% auditáveis" com fichas críticas completas resultará em rejeição automática, independentemente
% da qualidade técnica ou científica do artigo. Este é um padrão non-negotiable para todas
% as submissões neste ecossistema.